\section{Model Invocation and Empty-Output Diagnostics}
\label{app:model-invocation}

The main run used provider-compatible profiles. All models omitted \code{temperature}. DeepSeek and Kimi used no max-completion token cap in the main operational profile after capped profiles produced empty or malformed outputs in diagnostics. In the broader model-selection cap/no-cap diagnostic, the affected \code{120} cap empty-content calls ended with \code{finish_reason=length}; rerunning the same short JSON smoke test without \code{max_completion_tokens} returned parseable JSON.

The no-cap diagnostic did not indicate runaway visible output. In the opaque main run, \code{deepseek-v4-pro} averaged 370.340 visible output tokens. \code{kimi-k2.6} averaged 377.885 visible output tokens. Both had \code{overlength_rate=0.0}. Thus the empty-output issue is treated as an invocation-profile compatibility issue, not as evidence that the upstream model was unavailable or that the model required unusually long answers.

The model-availability and cap/no-cap diagnostics are recorded in the model-selection report and in the cap-versus-no-cap JSON artifact under \code{eha-mvp/results/}.
